% Integrated original source: EGAVaiello.pdf, pages 53--61.
% Dot: 17 archived batches. Two subsequent section-16 batches are collated separately.
\egasection{§14\quad 锥投影}

% SOURCE: unpacked/dot_batches/EGA第五卷_14锥投影构造_续译工程/EGA第五卷_14锥投影构造_校订续译.tex
\noindent\textbf{附记。}
我们已经在不同情境中使用过锥投影，尤其是在第8节末尾、10.4的表述以及其他地方。下面这组汇集的命题（“sorite”）无疑应提前放在本段的开头，也可能放在关于“Grassmann概形”的辅助段落中。

设$C=P(F)$是$P(E)=P$的一个线性子簇，它在$S$上的相对维数为$r-m-1$，即在$P$中的余维为$m+1$。于是$F$是$E$的一个秩为$r-m$的局部自由商模，$F=E/G$，其中$G$局部自由且秩为$m+1$。我们已用第二章的代数方法定义了态射
\[
p_c:P-C=P(E)-P(E/G)\longrightarrow P(G).
\]
我们将对它作几何解释，并依据下面的描述，将其称为以$C$为中心的\textbf{锥投影}。（附记：只假设$r-m-1$介于$-1$与$r-1$之间，即$m$介于$0$与$r$之间，并无其他限制。\cmt{V-H-L001}{本节的$r$表示$P(E)$的相对维数，因$\operatorname{rk}F=r-m$和$\operatorname{rk}G=m+1$，此处$\operatorname{rk}E=r+1$。第12节在第43页另以$r$表示$E$的秩，故本节记号已切换；不能把此处$r-m-1$静默改为$r-m-2$，也不能把整个$P(E)$的相对维数误读成$r-1$。这里$P(E)$采用局部自由秩一商模的约定；因此中心$P(E/G)$对应商模$E/G$。当$m=r$时，该商模为零，中心为空，射影维数$-1$是空射影空间的约定。保留原文的整个指标范围，不额外要求中心非空。}

为此，先将$P(G)$解释为
\[
\operatorname{Grass}^{m}(P)=\operatorname{Grass}_{r-m+1}(E)
\]
的一个闭子概形。\cmt{V-H-L002}{底本等式右端印为$\operatorname{Grass}_{r-m+1}(P)$，但随后的函子态射右端明确写$\operatorname{Grass}_{r-m+1}(E)$，并使用$E$的秩为$r-m+1$的商模。因此正文改为$E$；上标$m$表示射影线性子空间的余维，下标$r-m+1$表示局部自由商模的秩，不混用二者。}
这是由显然的函子同态
\[
P(G)\longrightarrow\operatorname{Grass}_{r-m+1}(E)
\]
得到的：对于$G$的每个可逆商模$G/G'$，取$E$的局部自由商模$E/G'$，其秩为$(r-m)+1$；在任意基变换之后也作同样的构造。）

上述函子同态是单态；又因源在$S$上固有、靶在$S$上分离，所以它是闭浸入。\cmt{V-H-L003}{底本重复写“The above homomorphism of functors is a homomorphism”，但仅说是同态不足以作出闭浸入结论。由商$E\twoheadrightarrow E/G'$的核可恢复$G'\subset G$及商$G/G'$；任意基变换后仍如此，故此函子态射是单态。源固有、靶分离使该态射固有，固有单态为闭浸入。这是本处将第二个“homomorphism”订为“单态”的理由，并不增补后续定理。}

更一般地，我们可能需要明确说明，如何将$G$的Grassmann概形，也就是$P(G)$的Grassmann概形（在函子的意义下），闭浸入到$E$的Grassmann概形，也就是$P(E)$的Grassmann概形中。所得态射的像（在函子的意义下）由$P$中具有所需维数、且包含$C$的线性子簇$L$组成。

在我们所考察的情形，即维数如上指定的情形，将这个像记为$Q(C)$。把$P(G)$与$Q(C)$等同后，锥投影态射
\[
p_c:P-C\longrightarrow Q(C)\subset\operatorname{Grass}^{m}(P)
\]
无非是把$P-C$的每个截面对应到$P$中唯一一个同时包含$C$和该截面、且余维为$m$的线性子簇$L$。（当然，所谓包含一个截面，是指该截面通过$L$分解。）

% RESTORED DIRECTLY FROM SOURCE, pages 54--55:
现在，若给定态射$f:X\to P$，便可考虑复合态射
\[
X-f^{-1}(C)\longrightarrow P-C\longrightarrow Q(C).
\]
我们可将它称为$X$的相对于$f$、以$C$为中心的锥投影，记为$p_c^X$，或简记为$p_c$。须指出，这个态射一般并不定义在整个$X$上；更确切地说，它定义在整个$X$上，当且仅当$f^{-1}(C)=\varnothing$，也就是说，$f(X)$与投影中心$C$不相交。\cmt{V-H-L004}{本段至下文“参看第10节和第12节”原未进入已恢复的Dot分批正文，本次据底本第54页末至第55页逐段补译。底本起句“If not, we have an $f:X\to P$”的连接词有误，本句按随后明确给出的构造译为“现在，若给定”；同段的$f(x)$按所说态射定义域应为像集$f(X)$。原空集符号$\phi$按其意义统一为$\varnothing$。}

我们将用前面各节中使用的构造，对这个态射给出另一种解释。为此，沿用别处已经引入的记号，考察
\[
\begin{tikzcd}[column sep=large]
X & X^{(m)}_{Q(C)}
   =X^{(m)}\times_{\operatorname{Grass}^{m}(P)}Q(C)
   \arrow[l,"q"'] \arrow[d,"p"] \\
  & Q(C).
\end{tikzcd}
\]
另一方面，注意到$q$诱导同构
\[
q':q^{-1}\bigl(X-f^{-1}(C)\bigr)
\xrightarrow{\ \sim\ }X-f^{-1}(C),
\]
而且立即可见，$p_c$正是$p'q'^{-1}$，其中$p'$表示$p$在
$q^{-1}\bigl(X-f^{-1}(C)\bigr)$上的限制。因此，利用$q'$直接作上述等同后，我们可以说，$p_c$就是态射$p$在
\[
X-f^{-1}(C)\subset X^{(m)}_{Q(C)}
\]
上的限制。\cmt{V-H-L005}{底本图式把纤维积的乘号排成字母$X$，基底简写为$\operatorname{Grass}^{m}$；正文按所述基变换写成$\times_{\operatorname{Grass}^{m}(P)}$。本段末尾的环境空间印作$X(M)^{p_c}_{Q(C)}$，与本段唯一已定义的关联空间不合；由$q'$所给开集等同可确定其应为$X^{(m)}_{Q(C)}$，正文据此订正，不把原排式另立为对象。}

因此，为方便起见，仍把上述态射记为$p_c^X$或$p_c$，并称之为$X$相对于$f:X\to P$、以$C$为中心的\emph{扩展锥投影}〔Mettons〕。\cmt{V-H-L006}{底本在此夹有“[(Mettons)]”，意为“姑且这样说／定名”；正文保留这一编写痕迹。这里“上述态射”指关联空间到$Q(C)$的$p$，不把扩展投影的定义域仍写成$X-f^{-1}(C)$。}
这样，受限锥投影的性质就归结为扩展锥投影的性质；后者已在别处得到系统研究，或者应当已经得到研究\footnote{底本脚注67：“Ask A.G.”，即“询问A.G.”。}；这也是有意义的（参看第10节和第12节）。

% SOURCE: unpacked/dot_batches/EGA第五卷_14一般中心_续译工程/EGA第五卷_14一般中心_校订续译.tex
主要的问题是：若$S=\operatorname{Spec}(k)$，并取$C$为$\operatorname{Grass}^{m+1}(P)$中的一般中心，\cmt{V-H-L007}{底本在这里附“[illegible, ask AG]”，译为：［不可辨，询问AG］。这条编者批注不提供可恢复的数学条件，不能据此猜补。原式的射影空间小写$p$统一为本节$P$。}那么$X$的锥投影具有什么性质？这要求作基变换$k\to k(\eta)$；也就是说，此时$C$确实是$P_{k(\eta)}$中的线性子簇。\cmt{V-H-L008}{底本先写$k\to K(\eta)$，随后使用$k(\eta)$，这里统一为一般点$\eta$的剩余域$k(\eta)$。底本本句写“a linear subvariety of $X_{k(\eta)}$”；但中心由$\operatorname{Grass}^{m+1}(P)$参数化，须位于$P$的基变换中，并不要求包含于$X$。第56页紧接着明确写$C_\eta\subset P_{k(\eta)}$，故正文据此改为$P_{k(\eta)}$。}

由已经重复过多次的标准论证，可以对上述Grassmann概形中某个非空开集的点所对应的锥投影，推出类似的性质；最后，若$k$是无限域，就可推出存在一个在$k$上定义的$C$，而当$m<r$时事实上有无穷多个这样的$C$。它们是$P$本身中的线性子簇，不必改变基域，并给出具有所讨论性质的锥投影。\cmt{V-H-L009}{底本写“since $k$ is infinite (if)”以及“a (in fact of an infinity of) $C$”。前者按其括注译为“若$k$是无限域”。后者须排除本节允许的$m=r$边界：此时中心为唯一的空线性子空间，参数Grassmann概形仅为一点，不存在无穷多个不同中心。$0\le m<r$时，该Grassmann概形维数$(m+1)(r-m)>0$，其标准仿射胞腔与任一非空开集相交为非空开集，在无限域上有无穷多个有理点，故正文补明这一范围。这里保留作者对“所讨论性质”的概括，不宣称每种性质都自动满足从一般点推广到开集的条件。}

宜将（或将来应将）这类一般性说明，同第4节和第7节中同类的说明汇集在一起；这些说明我们已经或多或少地隐含使用过，例如在第13节中。

此外，在这一联系下，也宜考察$X$上的一个层$F$的相对性质，并取它在$X^{(m)}_{Q(C)}$上的逆像$F^{(m)}_{Q(C)}$。在这里所述的具体情形中，还需要简化记号。我建议使用$X(C)$和$F(C)$，或者在不会混淆时，简记为$\widetilde X$和$\widetilde F$（注意：这里的$F$与本节开头的$F$不同）。\cmt{V-H-L010}{底本建议“$X(C)$ and $F(S)$”，但本句的对象是固定中心$C$所给关联空间及其上的逆像层，紧前均以$Q(C)$作下标，$S$则是既定基底而非这里变化的中心。因此正文统一为$F(C)$；底本$F(S)$原写法保留于此。这里的层$F$不与第54页定义中心$C=P(F)$的局部自由商模混同，原作者的警告完整保留。}

